% !PRINTED 602
\weberpage{602}
%<<CONT
we have called the discriminant of the number $\omega$. The ratio $D:d$ is an integral rational square number. The two roots of (1) receive the expression

\weq{3}{\omega = \frac{b+\sqrt{D}}{2c}, \quad \omega' = \frac{b-\sqrt{D}}{2c}.}
If we set
\weq{4}{D = e^2 d,}
then
\weq{5}{\omega = \frac{b+e\sqrt{d}}{2c}, \quad \omega' = \frac{b-e\sqrt{d}}{2c},}
and on account of (1) these are then and only then integers, when $c=1$.

Now $D\equiv 0$ or $\equiv 1\ (\mathrm{mod}\ 4)$, and if, therefore, $D$ is even, it is divisible by 4; $b$ is in this case even, and since $d$ cannot be divisible by 4, $e$ must also be even.

But if $D$ is odd, then $b$ and $e$ are also odd. This case, however, as (4) shows, can occur only when $d\equiv 1\ (\mathrm{mod}\ 4)$.

From this we obtain the following result:

1. \gesperrt{If $d\equiv 2$ or $\equiv 3\ (\mathrm{mod}\ 4)$, the number}
\[ \omega = x + y\sqrt{d} \]
\gesperrt{is then and only then an integer, when the rational numbers $x, y$ are integers.} (Quadratic irrational numbers of the first kind, after \gesperrt{Dirichlet}.)

2. \gesperrt{If $d\equiv 1\ (\mathrm{mod}\ 4)$, then}
\[ \omega = \frac{x+y\sqrt{d}}{2} \]
\gesperrt{is then and only then an integer, when the rational numbers $x,y$ are either two even or two odd integers.} (Quadratic irrational numbers of the first or second kind, according as $e$ is even or odd.)

We can also express the latter thus:

3. \gesperrt{If $d\equiv 1\ (\mathrm{mod}\ 4)$, then}
\[ \omega = x + y\,\frac{1+\sqrt{d}}{2} \]
\gesperrt{is then and only then an integer, when the rational numbers $x$ and $y$ are integers.}
